%!TEX root = main.tex
\section{Conclusion}
\label{sec:conclusion}

We presented SensitiveAI, a bilingual language-aware hybrid of lexicon rules
and multilingual transformers for Vietnamese and English hate speech detection.
The system combines a lightweight language detector, language-aware
tokenization, two independently editable rule engines (668 Vietnamese entries,
818 English entries), per-language mDeBERTa-v3-base checkpoints, and
language-aware score fusion with per-class thresholds.

The key results are:
\begin{itemize}
    \item On ViHSD, the hybrid is macro-F1 neutral (0.6575 $\to$ 0.6602,
          $p=0.29$), but its 38 changed decisions are fully decomposable into
          rule-driven and threshold-driven causes, with all 23 corrections
          rule-supported.
    \item On the 9{,}158-sample bilingual evaluation (full ViHSD $+$ Davidson splits), the pipeline
          reaches 87.10\% accuracy, macro-F1 0.7781, 9.07 ms latency and
          110.2 texts/s throughput (serialised run, $\alpha=0.70$, P95 12.69 ms).
    \item The bilingual comparison exposes a structural asymmetry: English
          excels on OFFENSIVE (F1 0.9270) while Vietnamese excels on HATE
          (F1 0.6051); the 818-entry English lexicon still has only 86 HATE
          terms versus 594 for Vietnamese.
    \item Fusion improves calibration: the full-pipeline ECE is 0.076 versus
          0.114 (VI backbone-only) and 0.102 (EN backbone-only), confirming that
          fused scores are better calibrated than raw neural outputs.
\end{itemize}

The broader lesson is methodological: a hybrid evaluated only by its aggregate
metric looks like a null result, whereas the same hybrid decomposed by channel
and direction yields an auditable decision procedure and a design agenda.
Future work includes evidence-gated fusion, balanced English HATE lexicon
expansion, a multilingual fallback checkpoint for code-mixed text, and a
human evaluation of whether the evidence record helps reviewers.